% Options for packages loaded elsewhere
\PassOptionsToPackage{unicode}{hyperref}
\PassOptionsToPackage{hyphens}{url}
%
\documentclass[
  11pt,
]{article}
\usepackage{lmodern}
\usepackage{amssymb,amsmath}
\usepackage{ifxetex,ifluatex}
\ifnum 0\ifxetex 1\fi\ifluatex 1\fi=0 % if pdftex
  \usepackage[T1]{fontenc}
  \usepackage[utf8]{inputenc}
  \usepackage{textcomp} % provide euro and other symbols
\else % if luatex or xetex
  \usepackage{unicode-math}
  \defaultfontfeatures{Scale=MatchLowercase}
  \defaultfontfeatures[\rmfamily]{Ligatures=TeX,Scale=1}
\fi
% Use upquote if available, for straight quotes in verbatim environments
\IfFileExists{upquote.sty}{\usepackage{upquote}}{}
\IfFileExists{microtype.sty}{% use microtype if available
  \usepackage[]{microtype}
  \UseMicrotypeSet[protrusion]{basicmath} % disable protrusion for tt fonts
}{}
\makeatletter
\@ifundefined{KOMAClassName}{% if non-KOMA class
  \IfFileExists{parskip.sty}{%
    \usepackage{parskip}
  }{% else
    \setlength{\parindent}{0pt}
    \setlength{\parskip}{6pt plus 2pt minus 1pt}}
}{% if KOMA class
  \KOMAoptions{parskip=half}}
\makeatother
\usepackage{xcolor}
\IfFileExists{xurl.sty}{\usepackage{xurl}}{} % add URL line breaks if available
\IfFileExists{bookmark.sty}{\usepackage{bookmark}}{\usepackage{hyperref}}
\hypersetup{
  hidelinks,
  pdfcreator={LaTeX via pandoc}}
\urlstyle{same} % disable monospaced font for URLs
\usepackage[margin=1in]{geometry}
\setlength{\emergencystretch}{3em} % prevent overfull lines
\providecommand{\tightlist}{%
  \setlength{\itemsep}{0pt}\setlength{\parskip}{0pt}}
\setcounter{secnumdepth}{-\maxdimen} % remove section numbering
\usepackage{mathrsfs}
\usepackage{esint}
\setmathfont{Asana Math}
\setlength{\emergencystretch}{3em}

\author{}
\date{23 September 2026}

\begin{document}

\hypertarget{harmonic-chern-representatives-at-the-realizable-einstein-calibration}{%
\section{Harmonic Chern representatives at the realizable Einstein
calibration}\label{harmonic-chern-representatives-at-the-realizable-einstein-calibration}}

23 September 2026. This application starts with a metric and an actual
conformally solvable Spin\(^c\) structure. No auxiliary form, scalar
band, spectral gap, or second-order budget is assumed. Its conclusion is
not another Einstein obstruction: the specified harmonic Chern
representative is quantitatively close to a bounded almost-parallel form
and to low connection-Laplacian modes.

\hypertarget{raw-hypotheses-and-statement}{%
\subsection{1. Raw hypotheses and
statement}\label{raw-hypotheses-and-statement}}

Let \((M,g)\) be a closed connected oriented smooth four-manifold,
normalized relative to the target \(\operatorname{Ric}=-g\). Fix a
Spin\(^c\) structure with real first Chern class \(a\), whose
unperturbed Seiberg--Witten equations are solvable for every metric in
\([g]\). The structure may depend on \(g\), but remains fixed when the
conformal weight varies. Write \[
 V=\operatorname{Vol}(g),\quad p=(a_g^+)^2>0,\quad
 r=\fint|\operatorname{Ric}+g|^2,\quad b=\frac{(-a^2)_+}{V},
 \quad w=\sqrt6|W^+|,
\] \[
 \delta=\frac32\int(s^2+2w^2)-72\pi^2p\ge0,
 \qquad \kappa=r+\delta/V.
 \tag{1}
\] Ricci uses tensor norms, forms use exterior-algebra norms, and Weyl
operators use Hilbert--Schmidt norms. The deficit equals
\(144\pi^2D_{\rm LB}\), or the parameter \(\beta=2\) in
\texttt{line1\_functional\_pencil.md}. Assume \(\kappa\le1/2880\).

Let \(H\) be the harmonic self-dual representative of \(a_g^+\), with
\(\|H\|_2^2=p\), and set \(\zeta=-2\pi H\). Our curvature convention is
\(\operatorname{proj}_{\mathcal H^+}(iF_A^+)=2\pi H\). Put
\(G(\kappa)=13\kappa+4\sqrt{2\kappa}\).

\textbf{Theorem.} For each \(0<\eta\le1/4\), there is a smooth actual
weighted-monopole form \(\omega\in\Omega^+(M)\) such that \[
 \|\omega\|_\infty\le2\sqrt3,
 \qquad \boxed{\fint|\omega-\zeta|^2\le49\kappa+\eta,}
 \tag{2}
\] \[
 \boxed{\fint|\nabla\omega|^2\le G(\kappa)+\eta,\qquad
 \fint|d\omega|^2\le\frac{13}{2}\kappa+\eta.}
 \tag{3}
\] Define \[
 A_{\omega,1/2}=\frac12
       \left(\frac{|\omega|^2}{3}I-\omega\otimes\omega\right),
\] \[
 d_\omega=\fint\left(|\operatorname{Ric}+g|^2
       +|W^+-A_{\omega,1/2}|^2+(|\omega|^2-2)^2\right),
 \quad B_\omega=\fint|\Delta_H\omega|^2.
\] The same smooth form satisfies, with universal constants, \[
 d_\omega\le720\kappa+\eta,\qquad
 B_\omega\le C(1+\kappa+b),\qquad
 (1+B_\omega)d_\omega\le C(1+\kappa+b)(\kappa+\eta).
 \tag{4}
\] Here \(d\omega\) is exterior differentiation and \(\Delta_H\) is the
nonnegative Hodge Laplacian. The first-order conclusions (2)--(3) do not
depend on \(b\); that cost enters only (4).

Let \(\mathscr L=\nabla^*\nabla\) on \(\Lambda^+\), let \(\mu_0\ge0\) be
its bottom eigenvalue, including zero, and let \(\Pi_{\le\Lambda}\)
project onto eigenvalues at most \(\Lambda>0\). Then the input harmonic
representative satisfies \[
 \boxed{
 \mu_0\le\frac{G(\kappa)}{2-\sqrt{720\kappa}},\qquad
 \fint|(I-\Pi_{\le\Lambda})\zeta|^2
       \le98\kappa+\frac{2G(\kappa)}{\Lambda}.}
 \tag{5}
\] No scalar or Hodge spectral gap, scalar-curvature sign, noncollapse,
Sobolev constant, or full-curvature bound is used. In particular (2)
does not follow from small exterior derivative via a uniform Hodge
inverse; its harmonic class comes from the actual curvature equation.

\hypertarget{producing-the-bounded-form-and-its-zeroth-order-budget}{%
\subsection{2. Producing the bounded form and its zeroth-order
budget}\label{producing-the-bounded-form-and-its-zeroth-order-budget}}

Put \(q=-s+w\), \(e=s+4\), \(z=s+2w\), and \(q_\pm=\max(\pm q,0)\). The
signed monopole estimate gives \(\Delta_+:=\int q_+^2-72\pi^2p\ge0\).
Completing the square, \[
 \delta=\Delta_++\int q_-^2+\tfrac12\int z^2,
 \qquad q-6=(z-3e)/2,\quad w-2=(z-e)/2.
 \tag{6}
\] Since \(e^2\le4|\operatorname{Ric}+g|^2\), \[
 \fint(q-6)^2\le\delta/V+18r\le18\kappa,
 \qquad \fint(w-2)^2\le\delta/V+2r\le2\kappa.
 \tag{7}
\] Set \(Q=\max(q,3)\) and \(f_0=\max\{\sqrt Q,(-s)_+/4\}\). Choose
smooth majorants \(f_j\ge f_0\) converging uniformly to \(f_0\).
Conformal solvability supplies \[
 D_A\Phi=0,\quad -iF_A^+=f_j\sigma(\Phi),\quad
 \psi=2\sqrt2\sigma(\Phi),\quad
 \rho=|\Phi|^2=|\psi|,\quad u=3\rho/2,
 \quad \omega=\sqrt3\psi/2.
 \tag{8}
\] With \(|\sigma(\Phi)|^2=|\Phi|^4/8\), the maximum principle gives
\(\rho\le4\) and \(u\le6\).

We record the moment inputs, including their truncation cost. Write \[
 \begin{split}
 D_0={}&2\int|d\rho|^2
 +\int(4\rho|\nabla_A\Phi|^2-|\nabla\psi|^2)\\
 &+2\int|d\psi|^2+2\int E_\lambda,\qquad
 E_\lambda=\langle W^+\psi,\psi\rangle+w\rho^2/3\ge0.
 \end{split}
\] All its summands are nonnegative, and \(27D_0/8=\int(qu^2-f_ju^3)\).
For \(C=\int Qu^2\), \(Z=\int\sqrt Q\,u^3\), \(E_Q=\int Q^2\), and
\(C_{\rm act}=\int f_j^2u^2\), \[
 C-Z\ge27D_0/8,\quad C\le E_Q^{1/3}Z^{2/3},\quad
 C_{\rm act}\ge72\pi^2p,\quad
 C_{\rm act}-C\le\tfrac{81}{4}\int e^2+o(1).
\] The last estimate uses \(Q^2-q_+^2\le9e^2\) and
\(f_0^2-Q\le9e^2/16\). Consequently, with
\(S=\Delta_++117\int e^2/4\le117\kappa V\), \[
 D_0\le4S/27+o(1),\quad
 \int Q(u-\sqrt Q)^2\le S+o(1),\quad
 0\le C_{\rm act}-72\pi^2p\le S+o(1).
 \tag{9}
\] These are the truncated Hölder moment estimates proved in the
functional-pencil note, Section 3, and Round 4, Section 4. All errors
here concern weight approximation on a fixed \((g,a)\); no limiting
spinor or common smoothing rate for a family is assumed.

The explicit zeroth-order bound follows without an eigenline choice: \[
 \fint(u^2-6)^2
 \le[234(1+2\sqrt3)^2+36]\kappa+o(1)
 \le4800\kappa+o(1),
 \tag{10}
\] because \(|u^2-Q|\le(1+2\sqrt3)\sqrt Q|u-\sqrt Q|\), and clipping
does not increase distance from six. Since \(|\omega|^2=u^2/3\), its
amplitude error is at most \(4800\kappa/9+o(1)\). The exact trace-free
matrix identity is \[
 |W^+-A_{\omega,1/2}|^2
 =\tfrac16(w-u^2/3)^2+\tfrac34 E_\lambda.
 \tag{11}
\] By (7), (9), and \(\int E_\lambda\le2S/27+o(1)\), its normalized
integral is at most \((2/3+4800/27+13/2)\kappa+o(1)\). Adding amplitude
and Ricci errors gives \(d_\omega\le720\kappa+o(1)\).

\hypertarget{differential-and-harmonic-information}{%
\subsection{3. Differential and harmonic
information}\label{differential-and-harmonic-information}}

The self-dual Weitzenböck formula and (9) give directly \[
 \begin{split}
 \int|\nabla\psi|^2
 &=2\int|d\psi|^2+2\int E_\lambda-\tfrac13\int z\rho^2\\
 &\le4S/27+(16/3)\sqrt{V\int z^2}+o(1).
 \end{split}
 \tag{12}
\] Multiplying by \(3/4\) proves (3). This controls the full covariant
derivative, without a Hessian estimate or Chern-square budget.

The actual curvature has harmonic projection
\(\operatorname{proj}_{\mathcal H^+}(f_j\psi)=-4\sqrt2\pi H\). Thus \[
 \|f_j\psi+4\sqrt2\pi H\|_2^2
 =\tfrac49(C_{\rm act}-72\pi^2p)\le4S/9+o(1).
 \tag{13}
\] Moreover \((f_0-\sqrt Q)^2\le f_0^2-Q\), so \[
 \fint|f_j-\sqrt6|^2
 \le\tfrac92r+
 \frac{2(\delta/V+18r)}{(\sqrt3+\sqrt6)^2}+o(1)
 \le9\kappa+o(1).
 \tag{14}
\] Use the identity \[
 2\sqrt2(\omega+2\pi H)
 =(\sqrt6-f_j)\psi+(f_j\psi+4\sqrt2\pi H)
\] and \(|\psi|\le4\). The squared triangle inequality gives \[
 \fint|\omega+2\pi H|^2
 \le4\fint|f_j-\sqrt6|^2+S/(9V)+o(1)
 \le49\kappa+o(1),
 \tag{15}
\] which identifies the prescribed harmonic Chern representative.

\hypertarget{second-order-budget-and-smoothing}{%
\subsection{4. Second-order budget and
smoothing}\label{second-order-budget-and-smoothing}}

Equations (6)--(7) bound the normalized squared integrals of
\(\operatorname{Ric},s,W^+\) by \(C(1+\kappa)\), and the cutoff bounds
\(\fint f_j^2\le C(1+\kappa)+o(1)\). Hence \[
 \int|F_A|^2=2\int|F_A^+|^2-4\pi^2a^2
 \le CV(1+\kappa+b)+o(1).
\] The positive-spinor connection curvature is controlled by
\(W^+,s,\operatorname{Ric}^{\circ},F_A\); \(W^-\) does not act on this
representation. Apply the inherited bounded-section estimate \[
 \|\nabla^2v\|_2^2\le2\|\nabla^*\nabla v\|_2^2
 +32\|v\|_\infty^2\|\operatorname{Ric}\|_2^2
 +33\|v\|_\infty^2\|\mathscr F\|_2^2,
 \quad \|\nabla v\|_4^2\le4\|v\|_\infty\|\nabla^2v\|_2
\] to \(\Phi\), whose rough Laplacian is controlled by
\((s+f_j\rho)\Phi\). Twice differentiating the parallel quadratic map
and using the self-dual Hodge formula gives
\(B_\omega\le C(1+\kappa+b)\). This is the second-order argument in the
functional-pencil note, Section 5, and Round 21, Section 2; it
differentiates neither the weight nor the Ricci error.

For each fixed metric and class, select one sufficiently large \(j\) so
that all normalized approximation errors meet their allotted fractions
of \(\eta\). All estimates hold for the same smooth actual solution. The
second-order bound holds for every sufficiently close majorant. This
proves (2)--(4), with no limiting-solution claim.

\hypertarget{connection-spectrum-and-smooth-approximate-projections}{%
\subsection{5. Connection spectrum and smooth approximate
projections}\label{connection-spectrum-and-smooth-approximate-projections}}

The amplitude bound implies
\(\fint|\omega|^2\ge2-\sqrt{720\kappa+\eta}>1\), which gives the
Rayleigh bound on \(\mu_0\). The spectral theorem gives \[
 \|(I-\Pi_{\le\Lambda})\omega\|_2^2
 \le\|\nabla\omega\|_2^2/\Lambda.
\] By (2)--(3), the normalized tail of \(\zeta\) is at most
\(98\kappa+2G(\kappa)/\Lambda+2\eta(1+1/\Lambda)\). The harmonic
representative and the spectral projection are fixed as
\(\eta\downarrow0\), proving (5) without convergence of spinors.

For a sequence with \(\kappa_j\to0\), even if \(b_j\to\infty\),
\(\fint|\zeta_j|^2=4\pi^2p_j/V_j\to2\). Choose \(\Lambda_j\to0\) with
\(G(\kappa_j)/\Lambda_j\to0\). Then all its normalized \(L^2\) mass lies
asymptotically in the low connection-Laplacian subspace. This does not
assert a pointwise bound on \(H_j\), nor is it a scalar-Laplacian
cluster statement.

The same form defines \(P_\omega=\omega\otimes\omega/2\). Algebra gives
\[
 \begin{aligned}
 \fint|W^++P_\omega-I/3|^2&\le d_\omega,\\
 \fint|P_\omega^2-P_\omega|^2&\le9d_\omega,\\
 \fint|\nabla P_\omega|^2&\le12(G(\kappa)+\eta).
 \end{aligned}
 \tag{16}
\] Outside \(\{|\omega|^2\le1\}\), of relative volume at most
\(d_\omega\), the endomorphism \(\omega\otimes\omega/|\omega|^2\) is a
rank-one projection with derivative energy at most
\(2V(G(\kappa)+\eta)\). This does not remove the exceptional set or give
an everywhere nonvanishing closed form. An \(L^2\) approximation also
does not permit substituting the potentially unbounded \(\zeta\) for
\(\omega\) in the quadratic Weyl model.

\hypertarget{an-actual-family-outside-every-fixed-negative-scalar-band}{%
\subsection{6. An actual family outside every fixed negative scalar
band}\label{an-actual-family-outside-every-fixed-negative-scalar-band}}

Take a closed negative Kähler--Einstein surface \((M,h)\), with
\(\operatorname{Ric}_h=-h\), carrying an actual canonical monopole
class; a product of two closed hyperbolic surfaces is an example. In
normal coordinates centered at \(x_0\), choose
\(\phi\in C_c^\infty(B_1(0))\) with \(\phi(0)=0\), \(d\phi(0)=0\), and
\(\Delta_{\mathrm{Eucl}}\phi(0)>0\), for the nonnegative-Laplacian sign
convention. Set \[
 f_\varepsilon(x)=\varepsilon\phi(x/\varepsilon),\qquad
 g_\varepsilon=e^{2f_\varepsilon}h,
\] extending the bump by zero outside its coordinate ball. In dimension
four, \[
 \int|\nabla_h^2f_\varepsilon|^2=O(\varepsilon^2),\quad
 \int|df_\varepsilon|^4=O(\varepsilon^4),\quad
 \|f_\varepsilon\|_\infty=O(\varepsilon).
\] The conformal Ricci formula, and cancellation of conformal factors in
the four-dimensional squared norm integral of a covariant two-tensor,
give
\(\int|\operatorname{Ric}_{g_\varepsilon}+g_\varepsilon|^2 =O(\varepsilon^2)\).
Volumes converge to \(\operatorname{Vol}(h)\). Both \(p\) and
\(\int|W^+|^2\) are conformally invariant. Since \[
 s_{g_\varepsilon}=e^{-2f_\varepsilon}
       (-4+6\Delta_hf_\varepsilon-6|df_\varepsilon|^2),
\] integration of \(\Delta_h f_\varepsilon\) gives \[
 \int s_{g_\varepsilon}^2-\int s_h^2
 =36\int(\Delta_hf_\varepsilon-|df_\varepsilon|^2)^2
       +48\int|df_\varepsilon|^2=O(\varepsilon^2).
\] Thus \(\delta_\varepsilon=O(\varepsilon^2)\) and
\(\kappa_\varepsilon\to0\), while \[
 s_{g_\varepsilon}(x_0)
 =-4+6\varepsilon^{-1}\Delta_{\mathrm{Eucl}}\phi(0)
 \longrightarrow+\infty.
\] The application therefore genuinely includes metrics outside every
fixed negative scalar band. These are integral-error examples, not
pointwise Ricci-pinched Yau26 inputs, and not new smooth topologies.

\hypertarget{increment-and-scope}{%
\subsection{7. Increment and scope}\label{increment-and-scope}}

The supplied form lies at the allowed Round-21 parameter \(t=1/2\),
where spectral non-saturation gives no contradiction. The moment and
Hessian mechanisms are inherited. Their all-scalar composition removes
the scalar band from the earlier actual-form supply. No scalar spectral
gap is needed when integral Ricci error is an input; this does not
remove the spectral-gap term from the different theorem deriving Ricci
control from period data alone. Identity (12) also removes the
negative-square cost from its first-order and low-mode conclusions. The
useful added output is the canonically identified harmonic alignment and
spectral/projection information, not a new topological obstruction or
existence theorem.

At zero defect this recovers the realizable negative Kähler--Einstein
equality case. At positive defect it supplies quantitative integral
structure, not global Kähler integrability or geometric compactness.

\hypertarget{proof-sources-and-verification}{%
\subsubsection{Proof sources and
verification}\label{proof-sources-and-verification}}

\begin{itemize}
\tightlist
\item
  \texttt{line1\_functional\_pencil.md}, Sections 2--5: actual
  all-scalar form production, specialized to parameter two.
\item
  \texttt{round4\_rigidity.md}, Sections 2--4: cutoff and moment
  estimates.
\item
  \texttt{round21\_analytic.md}, Section 2: the bundle Hessian estimate.
\item
  \texttt{line1\_scalar\_band\_remainder.md}, equations (9)--(10):
  earlier scalar-band differential and weighted harmonic remainders.
\item
  \texttt{line1\_composed\_application\_check.py}: exact finite
  coefficients and tensor identities, not certification of gauge
  solvability or PDE.
\end{itemize}

\end{document}
